\begin{afterword}
\summaryhead

The post offers an evolutionary explanation for why explanations fail. Humans evolved in small hunter-gatherer bands without writing, where, it argues, everyone shared the same background knowledge. Any new fact was at most one ``inferential step'' from what everyone already knew. In such a world, a person who asserted something without obvious support was a liar or a fool, and a person who failed to see something obvious was stupid or obstinate.

Modern knowledge is different: scientific conclusions can be a hundred steps from common knowledge. The author argues that we still expect short distances, and that this, together with the illusion of transparency, explains why scientists struggle to talk to lay audiences and to each other. His example is a biologist who defends evolution as ``the simplest explanation'' to a listener who does not share the scientist's respect for simplicity.

The advice is to build an argument step by step from what the audience already accepts, and not to let them see that you think you are far ahead of them.

\responsehead

The practical advice in this post is ordinary: start from what your audience knows. It takes two sentences. The rest of the post is an evolutionary story told to explain a finding that already had a name.

The story has no evidence behind it. The one source the post supplies, a link on the words ``200 people,'' goes to a page about Dunbar's number, which is 150 and measures social relationships, not band size. That page gives hunter-gatherer bands of 30 to 50 inside tribes of 500 to 2,500. The claim that ``all background knowledge is universal knowledge'' is asserted with ``period'' in place of an argument. It also forgets children. Every band has always had members many steps behind the adults, and teaching them is the long explanation every human group has had to give. The post's premise erases an ancestral activity that bears directly on its subject.

The story also explains nothing that its own sources do not already explain. The post names the illusion of transparency and self-anchoring and then adds the ancestral environment as a third cause. No observation in the post could tell the three apart. The effect itself, that people who know something cannot easily imagine not knowing it, was named ``the curse of knowledge'' by Camerer, Loewenstein and Weber in 1989. The post does not mention it.

The advice that the post does give concerns presentation. It tells the reader not to let the audience see that they think themselves a dozen steps ahead; it never asks whether they are. The audience's sense of being talked down to is treated as a misreading.

The post then breaks its own rule. It warns against being seen to give an argument more weight than the audience grants it, and in the next sentence tells the reader that simplicity ``is'' a knockdown argument for evolution. It gives no reason, and it describes the people who agree as ``raised to revere Occam's Razor.''

In short: an uncredited finding wrapped in an untested origin story, with advice about presentation in place of a method for explaining.
\end{afterword}
